% Paleta común para todos los diagramas
\colorlet{cSensor}{blue!60!black}
\colorlet{cProc}{orange!80!black}
\colorlet{cAct}{green!45!black}
\colorlet{cApp}{violet!70!black}
\colorlet{cPower}{gray!70!black}

% Colores de los gráficos de datos (paleta categórica validada para daltonismo)
\definecolor{serie1}{HTML}{2A78D6}
\definecolor{serie2}{HTML}{EB6834}
\definecolor{serie3}{HTML}{1BAF7A}
\definecolor{tinta}{HTML}{52514E}

\tikzset{
  >={Stealth[length=2mm]},
  flecha/.style={->, thick, draw=black!75},
  caja/.style={rectangle, rounded corners=2pt, draw=#1, fill=#1!10, thick,
    minimum width=2.3cm, minimum height=0.85cm, align=center, inner sep=2pt},
  etiqueta/.style={font=\scriptsize, text=black!70, fill=white, inner sep=1pt},
}

% Estilo base de los gráficos. Las diapositivas cambian el tamaño con \renewcommand.
\newcommand{\anchografico}{9cm}
\newcommand{\altografico}{5.2cm}
\pgfplotsset{
  compat=1.18,
  grafico/.style={
    axis lines=left,
    axis line style={draw=black!45},
    tick style={draw=black!45},
    ymajorgrids, grid style={draw=black!12},
    tick label style={font=\scriptsize, text=tinta},
    label style={font=\footnotesize, text=tinta},
    legend style={draw=none, fill=none, font=\scriptsize, text=tinta},
    legend cell align=left,
    every axis plot/.append style={line width=0.8pt},
    /pgf/number format/use comma,
  },
}

% Mientras los archivos de figuras/datos/ tengan valores de ejemplo, cada
% gráfico lleva la marca "Datos de ejemplo". Al cargar las mediciones reales,
% cambiar a \datosdeejemplofalse en main.tex y en slides.tex.
\newif\ifdatosdeejemplo
\datosdeejemplotrue
\newcommand{\marcaejemplo}{%
  \ifdatosdeejemplo
    \node[anchor=north east, font=\scriptsize\bfseries, text=red!70!black,
      fill=white, draw=red!70!black, rounded corners=1pt, inner sep=2pt]
      at (rel axis cs:0.98,0.98) {Datos de ejemplo};
  \fi}

% Dato que falta medir o confirmar. Sale en rojo para que no pase desapercibido.
\newcommand{\pendiente}[1]{\textcolor{red!70!black}{[#1]}}

% Foto o captura: \imagen{archivo}{ancho}{alto}{qué debe mostrar}
% Si Imagenes/archivo existe, se inserta. Si no, se dibuja un recuadro con la
% descripción, así el documento compila mientras faltan las fotos.
\newcommand{\imagen}[4]{%
  \IfFileExists{Imagenes/#1}{%
    \includegraphics[width=#2, height=#3, keepaspectratio]{Imagenes/#1}%
  }{%
    \begin{tikzpicture}
      \node[draw=black!35, dashed, fill=black!4, rounded corners=2pt,
        minimum width=#2, minimum height=#3, text width=0.85\dimexpr#2\relax,
        align=center, font=\scriptsize, text=black!60] {%
        \textbf{Imagenes/#1}\\[1mm]#4};
    \end{tikzpicture}%
  }}
